% Author theory draft for R126. No empirical result is asserted by this file.
\section{Computing the decision content of paid Bellman information}
\label{sec:computed126}

The complete-support characterization identifies the information in a
Bellman calculation, but an unverified convex-program output does not yet
give an economic certificate. We now give a finite calculation which can
stop at any resource limit with a valid original-action bound. It combines
tangents taken at different points, verifies their aggregate by rational
arithmetic, and retains the stronger of this calculation and the already
available conventional bounds. Cutting-plane optimization is classical
\citep{Kelley1960Cuts126}; smooth-convex interpolation and information-based
worst-case calculations also have an established literature
\citep{THG2017Interpolation122,FloreaNesterov2024Oracle125}. The contribution
here is the explicit passage from the original, globally valid Bellman
support record to a finitely verified action and regret calculation. Its
economic information and certification rights are identical for a neural
critic, a matched ordinary hinge, a polynomial, a spline, or an online
solver supplied with that record.

\subsection{A finite certificate from numerical cut proposals}

Fix the state, risk law, feasible interval, original optimal-future
objective, and paid record of Section~\ref{sec:fullsupport125}. Write its
nonempty finite-dimensional compatibility set as
\[
 P=\{z\in\mathcal B:c_i(z)\le0,
                \ i=1,\ldots,m\},\qquad
 \mathcal B=\prod_k[\underline z_k,\overline z_k].
\]
The box is finite and rational. Linear restrictions, including equalities
written in both directions, are among the $c_i$. For a requested original
contrast $q(b)-q(a)$, let $H(z)=h_0+h^\top z=q(a)-q(b)$.
At arbitrary rational points $v_r$, choose valid rational subgradients
$s_r\in\partial c_{i_r}(v_r)$ and store the exact affine cuts
\[
 A_r=s_r^\top,\qquad b_r=s_r^\top v_r-c_{i_r}(v_r),
 \qquad A_rz\le b_r.
\]
The points need not be feasible, need not coincide, and need not be the
solutions of an optimization problem. All nonlinear support constraints
have explicit rational subgradients; the supplement gives the formulas.

\begin{theorem}[A finite verified bound from a shared cut record]
\label{thm:cuts126}
For any nonnegative rational weights $y_r$, define
\begin{equation}\label{eq:cutbound126}
 R=h+A^\top y,\qquad
 B=h_0-y^\top b+
       \sum_k\min\{R_k\underline z_k,R_k\overline z_k\}.
\end{equation}
Then every compatible original objective satisfies
$q(b)-q(a)\le-B$. Neither primal feasibility, stationarity, strong
duality, nor a successful solver exit is required. A numerical linear
program may propose the weights; sign checks, rationalization, reconstruction
of each exact cut, and evaluation of \eqref{eq:cutbound126} determine
whether a certificate is accepted.

If a rational point $z^{\rm f}\in P$ is independently verified, then
\begin{equation}\label{eq:cutgap126}
 -H(z^{\rm f})\le U_{\rm full}(a,b)\le-B,
 \qquad
 0\le -B-U_{\rm full}(a,b)\le H(z^{\rm f})-B.
\end{equation}
Thus a checked gap can certify accuracy relative to the exact declared
information bound. Without such a gap, the returned upper bound is valid
but is not certified to be the sharp information bound.

Intersecting each accepted bound with all earlier valid bounds for the
same original contrast gives a nonincreasing sequence of upper bounds.
The cut record may be reused for several linear contrasts on one common
node set, provided all arithmetic and optimization work is charged.
\end{theorem}

The zero-weight choice supplies a valid, possibly weak box bound. A failed
optimization, an unhelpful rationalized dual, or the imposed cut limit
therefore cannot invalidate the incumbent certificate. Feasibility of
the original information class is a scientific premise established by
the original calculation; it is not inferred from a numerical solver's
message. If an exact constraint check detects inconsistency in a supplied
record, the record is rejected rather than used to prove a favorable
decision.

The finite computation differs from an asymptotic convergence claim.
With exact linear subproblem solutions, bounded subgradients, and a cut
for every convex constraint at every current minimizer, the nested outer
relaxations converge in value to the exact compatibility program. A proof
is supplied in the supplement. A capped floating-point implementation
uses only \eqref{eq:cutbound126}; it does not inherit that limiting
statement as a finite accuracy certificate.

The information boundary also governs comparisons. The cheap conventional
calculation clips transported derivative envelopes by the same original
cap $[-M,M]$ and uses every available original-optimum lower bound
in its endpoint enclosure. If such an optimum bound has not been encoded
as a constraint of $P$, intersecting with it remains valid for the
actual original objective, but is not a tighter numerical solution of
$U_{\rm full}$. The raw cut bound and the final original-policy
intersection are reported separately. No information omitted from a
comparator is attributed to the convex solver or to neural architecture.

\subsection{From pairwise contrasts to original-policy regret}

The same record can answer a stronger economic question: how much
original value could still be lost by returning a particular action?
Define its information-compatible regret by
\begin{equation}\label{eq:inforegret126}
 \Gamma(a)=\sup_{q\in\mathcal Q_{\rm full}}
               \{q(a)-\min_{x\in[0,A]}q(x)\}.
\end{equation}
It is a property of the declared original information. It is not a
fitting loss, a posterior mean, or a claim about a different recursive
preference. Introduce an action grid
$0=x_0<\cdots<x_n=A$ containing $a$, with maximum spacing $h$.
These are additional unknown interpolation nodes, not newly purchased
economic observations. Theorem~\ref{thm:fullinterpolation125} ensures
that the augmented finite program represents the same function class.

\begin{theorem}[Finite decision resolution under unchanged information]
\label{thm:gridregret126}
Put $u_j=U_{\rm full}(x_j,a)$. Then
\begin{equation}\label{eq:gridsandwich126}
 \max_j u_j\le\Gamma(a)
        \le\max_j u_j+\frac{Lh^2}{8}.
\end{equation}
The constant $1/8$ is sharp, even for a compatible class consisting of
one rational quadratic. Both endpoints must be in the grid; a constrained
boundary optimum is then represented exactly.

More generally, let $\widehat u_j\ge u_j$ be any verified upper
contrasts, setting $\widehat u_j=0$ at $x_j=a$. For a cell of length
$d_j=x_{j+1}-x_j$, put
\begin{align}
 G_j(a)&=\max_{0\le t\le1}
   \left\{(1-t)\widehat u_j+t\widehat u_{j+1}
          +\frac{Ld_j^2}{2}t(1-t)\right\},\nonumber\\
 G(a)&=\max_jG_j(a).
 \label{eq:cellregret126}
\end{align}
Then $0\le\Gamma(a)\le G(a)$. For rational inputs $G(a)$ is
an exact rational number: with $k_j=Ld_j^2/2$, the maximizing
argument is
\begin{equation}\label{eq:cellmax126}
 t_j^*={\rm proj}_{[0,1]}
       \frac{\widehat u_{j+1}-\widehat u_j+k_j}{2k_j}.
\end{equation}
If every $\widehat u_j-u_j\le\epsilon$, then
\begin{equation}\label{eq:gridaccuracy126}
 0\le G(a)-\Gamma(a)\le \epsilon+\frac{Lh^2}{8}.
\end{equation}
Every quantity here concerns the same original objective and the same
paid record. Refining the computational grid requires additional
certificate work, but no new risk or continuation-value evaluation.
\end{theorem}

For the unchanged implemented future $J(a)=q(a)+\eta(a)$,
$0\le\eta(a)\le\delta$, the returned action consequently has
original-policy regret at most $G(a)+\delta$. Intersect this with
every earlier valid regret bound. An absolute original-value interval
at $a$ is still required separately: grid resolution cannot identify a
free additive value level. If a policy is changed, a strictly negative
upper bound on the actual contrast $J(b)-J(a)$ is still required.
A reduction in worst compatible regret by itself does not prove a
strict gain for the realized economy.

There is also a finite robust-action implication. The function
$\Gamma$ is $M$-Lipschitz and $\mu$-strongly convex. If the same grid
is used for candidate actions, and the verified pairwise accuracy is
at most $\epsilon$ for every candidate, a minimizer $\widehat a$
of the computed $G$ satisfies
\begin{equation}\label{eq:robustaction126}
 \Gamma(\widehat a)-\min_{a\in[0,A]}\Gamma(a)
       \le\frac{Mh}{2}+\frac{Lh^2}{8}+\epsilon.
\end{equation}
This is an information-compatible minimax regret guarantee, not
identification of the actual optimal action and not an advantage due
to a neural encoding. The supplement states its proof and scope.

Finally, suppose any further computation for this one decision must
first cost at least $c_*>0$ resource units, priced at $\lambda$ units
of original objective, and exclude reuse for other decisions. Stopping
is weakly optimal against all such continuations whenever
\[
 \lambda c_*\ge
 \min\{K_{\rm old},G(a)+\delta\},
\]
where $K_{\rm old}$ is a previously valid original-policy regret
bound. This is a dominance conclusion conditional on an actual
compulsory cost and the declared original objective. It does not turn
observed host clocks into a deployment probability or an estimated
economic price.
